% =====================================================================
%  Bachelor's thesis -- Context-aware adaptive beaconing and multi-hop
%  awareness for infrastructure-free maritime proximity detection.
%
%  Compile with:  pdflatex main.tex  (run twice for the ToC / references)
%  Figures live in ./figures and are produced by make_figures.py from the
%  averaged simulation CSVs in metrics/averaged/.
% =====================================================================
\documentclass[12pt,a4paper,oneside]{report}

\usepackage[utf8]{inputenc}
\usepackage[T1]{fontenc}
\usepackage[a4paper,margin=2.7cm]{geometry}
\usepackage{graphicx}
\usepackage{amsmath,amssymb}
\usepackage{booktabs}
\usepackage{array}
\usepackage{xcolor}
\usepackage{caption}
\usepackage{subcaption}
\usepackage{listings}
\usepackage{algorithm}
\usepackage{algpseudocode}
\usepackage{enumitem}
\usepackage{tabularx}
\usepackage{siunitx}
\usepackage[hidelinks]{hyperref}

\definecolor{codebg}{rgb}{0.97,0.97,0.97}
\definecolor{codekw}{rgb}{0.10,0.30,0.65}
\definecolor{codecm}{rgb}{0.30,0.55,0.30}
\lstset{
  basicstyle=\ttfamily\small,
  backgroundcolor=\color{codebg},
  keywordstyle=\color{codekw}\bfseries,
  commentstyle=\color{codecm}\itshape,
  stringstyle=\color{purple},
  numbers=left, numberstyle=\tiny\color{gray},
  breaklines=true, frame=single, framesep=4pt,
  showstringspaces=false, captionpos=b, language=Python,
}

\newcommand{\bimin}{BI_{\min}}
\newcommand{\bimax}{BI_{\max}}
\newcommand{\nthr}{N_{\mathrm{thr}}}
\renewcommand{\arraystretch}{1.15}

\title{\textbf{Context-Aware Adaptive Beaconing and Multi-Hop Awareness\\for Infrastructure-Free Maritime Proximity Detection}}
\author{Candidate Name\\\small ID number 0000000}
\date{Academic Year 2024/2025}

\begin{document}

% ---------------------------------------------------------------- title
\begin{titlepage}
\centering
{\scshape\LARGE Sapienza University of Rome \par}
\vspace{0.4cm}
{\scshape\large Faculty of Information Engineering, Informatics and Statistics \par}
{\large Bachelor's Degree in Computer Science \par}
\vfill
{\huge\bfseries Context-Aware Adaptive Beaconing and\\[0.3em] Multi-Hop Awareness for Infrastructure-Free\\[0.3em] Maritime Proximity Detection \par}
\vfill
\begin{flushleft}
\large
\textbf{Candidate}\\
Candidate Name --- ID number 0000000\\[1.2em]
\textbf{Advisor}\\
Prof.\ Gaia Maselli\\
\end{flushleft}
\vfill
{\large Academic Year 2024/2025 \par}
\end{titlepage}

% ---------------------------------------------------------------- abstract
\chapter*{Abstract}
\addcontentsline{toc}{chapter}{Abstract}

Swimmers, snorkelers, kayakers and divers are largely invisible to the larger
vessels that share coastal waters: the safety technologies designed for big
ships (AIS, satellite distress beacons, cellular tracking) are too expensive,
too slow, or dependent on infrastructure that does not exist offshore. A
promising alternative is to turn the Wi-Fi radio that every modern device
already carries into an \emph{infrastructure-free proximity beacon}: smart
buoys periodically broadcast short, location-aware frames that nearby vessels
listen for. The open question is how such broadcast beaconing should behave so
that it stays reliable, energy-frugal and \emph{aware} of the wider
neighbourhood as the density of users changes.

This thesis extends an event-driven simulator for IEEE~802.11 broadcast
beaconing and uses it to study two design axes that the existing literature
treats only in isolation: \emph{how often} a buoy should beacon, and \emph{how
far} the awareness it produces should propagate. Along the first axis we
introduce and evaluate \textbf{ACAB} (Adaptive Context-Aware Beaconing), a
scheduler that blends three locally observable signals --- neighbour density,
time since last contact, and node mobility --- into a single beacon-interval
decision, and we compare it against fixed-interval static beaconing (SBP) and
density-only adaptive beaconing (ADAB). Along the second axis we implement and
contrast two awareness-propagation mechanisms: \textbf{append} mode, which
piggybacks a hop-bounded neighbour list onto every beacon, and
\textbf{forwarded} mode, which actively relays beacons under a novel
density-aware forwarding gate that re-uses the scheduler's own back-off
intensity to suppress redundant relays. The channel model is calibrated from
over-water field measurements, mobility follows a Random-Waypoint model, and
the population churns over time to mimic users entering and leaving the water.

Across $30$ averaged runs per configuration, three buoy densities-per-interval
and node counts from $20$ to $100$, the experiments establish three main
results. First, \emph{context-aware adaptation pays for itself}: ACAB holds the
collision rate below $0.5\%$ where static beaconing climbs past $14\%$, and it
does so while transmitting up to $93\%$ fewer beacons than static and $62\%$
fewer than ADAB, with \emph{no} loss --- indeed a small gain --- in
per-receiver delivery ratio. Second, \emph{per-link reliability and
network-level awareness must be measured separately}: under a realistic
path-loss channel the per-receiver delivery ratio is bounded by propagation,
not contention, so the true value of adaptation surfaces in collision rate,
channel airtime and energy rather than in raw delivery. Third, and most
surprisingly, \emph{cheap passive awareness beats expensive active flooding}:
append mode discovers an equal or larger fraction of the network than forwarded
mode while generating three-to-five times less traffic and up to an order of
magnitude lower discovery latency. Together these findings argue that simple,
purely local information --- and parsimony in using the channel --- is
sufficient to build a scalable, infrastructure-free maritime safety beacon.

\clearpage
\tableofcontents
\listoffigures
\listoftables

% ================================================================
\chapter{Introduction}
\label{ch:intro}

\section{Context and motivation}
Coastal and near-shore waters concentrate two populations that are very hard to
make mutually visible. On one side there are large craft --- motorboats, ferries,
fishing vessels --- equipped with radar and the Automatic Identification System
(AIS)~\cite{imo-ais}. On the other side there are people in the water and the
small craft around them: swimmers, snorkelers, free-divers, paddlers and
kayakers. The instruments that protect the first population were never designed
to detect the second. AIS transponders are expensive, power-hungry and
optional for small vessels; satellite distress beacons such as EPIRBs and
PLBs~\cite{noaa-beacons} are meant for emergencies, not for the continuous,
second-by-second awareness that collision avoidance requires; and smartphone
tracking depends on cellular coverage that is unreliable a few hundred metres
from shore. The result is a persistent \emph{proximity-detection gap}: the most
vulnerable maritime users remain electronically silent to the craft most likely
to injure them.

Closing this gap calls for a technology that is simultaneously cheap,
energy-efficient, short-range and free of any fixed infrastructure. Wi-Fi
(IEEE~802.11~\cite{ieee80211}) is an attractive candidate precisely because it
is already everywhere. A buoy built around a commodity Wi-Fi chipset can
broadcast a small frame announcing ``I am here, at this position'' without
associating to any access point, without an Internet connection, and without
any central coordinator. A vessel needs only to listen. The hardware is a few
tens of euros, the range over open water is on the order of a hundred metres,
and the whole system degrades gracefully because there is nothing centralized
to fail.

\section{The scalability and awareness problem}
Wi-Fi was not, however, designed for safety-critical broadcast. Its Distributed
Coordination Function (DCF)~\cite{sharetechnote-dcf} arbitrates the medium with
a fixed set of CSMA/CA parameters~\cite{csma-ca} that do not adapt to how many
nodes are contending. Worse, \emph{broadcast} frames are never acknowledged: a
sender gets no feedback, cannot detect a collision, and cannot fall back on the
binary exponential back-off that rescues unicast traffic. As the number of
buoys in radio range grows, the probability that two of them transmit at once
grows with it, and beacons are silently lost. For a safety system this is the
worst possible failure mode --- a lost beacon is an undetected swimmer.

There is a second, less-discussed problem. A single beacon only tells a listener
about its \emph{immediate} sender. Yet a vessel approaching a group of swimmers
benefits enormously from knowing about the whole group, including members
temporarily hidden behind a wave or another hull. Awareness, in other words,
should propagate beyond one radio hop --- but every extra transmission spent
propagating it competes for the same scarce, unacknowledged channel that the
scalability problem is already straining. Beaconing frequency (\emph{how often}
to talk) and awareness propagation (\emph{how far} to spread what is heard) are
therefore not independent knobs: they draw on the same budget and must be
designed together.

\section{Contributions}
This thesis treats the two problems jointly. Building on an existing
discrete-event simulator for infrastructure-free Wi-Fi beaconing, it makes the
following contributions.

\begin{enumerate}[leftmargin=1.4em]
  \item \textbf{A context-aware adaptive scheduler (ACAB).} We design and
  implement a beaconing scheduler that fuses three locally observable signals
  --- neighbour density, recency of the last received beacon, and the node's own
  speed --- into one adaptive beacon interval, and we position it against a
  fixed-interval baseline (SBP) and a density-only adaptive scheme (ADAB).
  \item \textbf{Two multi-hop awareness mechanisms and a forwarding gate.} We
  implement \emph{append} mode (passive piggybacking of a hop-bounded neighbour
  list) and \emph{forwarded} mode (active relaying with a time-to-live), and we
  introduce a density-aware forwarding gate that re-uses the scheduler's own
  back-off intensity to suppress redundant relays before they ever contend.
  \item \textbf{A richer, decoupled metric suite.} We separate
  \emph{per-receiver delivery ratio} from \emph{network-level delivery ratio},
  add \emph{network-discovery percentage} and \emph{reaction latency}, and use
  them to show that the benefit of adaptation lives in channel load and
  collisions rather than in raw per-link delivery under a path-loss channel.
  \item \textbf{A grounded experimental study.} Using a field-calibrated
  channel, Random-Waypoint mobility and a churning population, we run a large
  averaged campaign ($30$ repetitions per point) and report results that are,
  in places, counter-intuitive --- most notably that cheap passive awareness
  outperforms expensive active flooding.
\end{enumerate}

\section{Research questions}
The work is organized around four questions.
\begin{description}[leftmargin=2em,style=nextline]
  \item[RQ1] How does node density affect the reliability and the channel cost
  of broadcast beaconing in an infrastructure-free network?
  \item[RQ2] Can a context-aware scheduler that uses only local information
  reduce channel load and collisions without sacrificing reliability, and how
  does it compare to a density-only scheme?
  \item[RQ3] What is the right way to propagate awareness beyond one hop ---
  passive piggybacking or active relaying --- and at what cost?
  \item[RQ4] Which performance metrics actually expose the value of adaptation
  once a realistic path-loss channel is in play?
\end{description}

\section{Structure of the thesis}
Chapter~\ref{ch:background} surveys maritime safety technology, the relevant
802.11 MAC mechanics, beaconing strategies and the role of custom simulation.
Chapter~\ref{ch:system} describes the ProSafe system and specifies the three
schedulers (SBP, ADAB, ACAB) and the two awareness mechanisms together with the
forwarding gate. Chapter~\ref{ch:simulator} details the simulator: its
event-driven core, the buoy state machine, the Random-Waypoint mobility model,
the calibrated channel, the interleaved transmission pipeline and the metric
suite. Chapter~\ref{ch:results} presents and discusses the experimental
results. Chapter~\ref{ch:conclusions} concludes and outlines future work.

% ================================================================
\chapter{Background and Related Work}
\label{ch:background}

This chapter provides the background needed to understand the problem and the
proposed solutions. It reviews why existing maritime safety technologies do not
fill the proximity-detection gap, recalls the parts of the IEEE~802.11 MAC layer
that govern broadcast behaviour, surveys beaconing and awareness-propagation
strategies, and motivates the use of a purpose-built simulator.

\section{Maritime and coastal safety technologies}
\subsection{Automatic Identification System (AIS)}
AIS is a VHF transponder system mandated by the International Maritime
Organization for ships above $300$ gross tonnage and for all passenger
ships~\cite{imo-ais}. It periodically broadcasts a vessel's identity, position,
course and speed so that equipped ships and shore stations can track one
another. AIS is excellent for ship-to-ship awareness in open water, but for
coastal proximity detection it has decisive drawbacks: transponders are
expensive and power-hungry, they are not carried by recreational users, their
use is licensed in many jurisdictions, and their update period (seconds to tens
of seconds for moving vessels) targets large-vessel navigation rather than
immediate short-range awareness.

\subsection{Distress beacons and cellular tracking}
EPIRBs and PLBs~\cite{noaa-beacons} transmit on $406$\,MHz to the
Cospas--Sarsat satellite constellation when manually activated. They are
life-saving in a genuine emergency but are unsuitable for routine awareness: a
false activation is costly, the satellite relay adds minutes of delay, and the
signal reaches rescue coordination centres rather than the nearby craft that
could actually avoid a collision. Smartphone applications backed by GPS and
cellular data can in principle warn of proximity, but they inherit the
unreliability of cellular coverage offshore, depend on infrastructure and a
data plan, and drain the battery quickly.

\subsection{The gap}
Every one of these technologies is either too expensive, infrastructure-bound,
or aimed at a different problem than continuous short-range detection. The
people most at risk --- those actually in the water --- stay invisible. The gap
calls for a low-cost, infrastructure-free, energy-aware solution operating at
ranges of tens to a couple of hundred metres.

\section{Wi-Fi as an infrastructure-free beacon}
Wi-Fi devices can either associate with an access point (infrastructure mode) or
talk peer-to-peer (ad-hoc / IBSS mode~\cite{ipcisco-ibss}). For proximity
detection neither is ideal as such; what matters is that a device can transmit
\emph{broadcast} frames that any in-range receiver decodes without association.
Re-purposing this capability, a buoy can emit a custom beacon carrying a device
identifier, a GPS position, a role flag, a battery level and a short list of
recently heard neighbours. The advantages are concrete: no access point is
needed, the hardware is ubiquitous and cheap, and the $100$--$300$\,m range
typical of $2.4$\,GHz operation matches the coastal use case. The challenges are
equally concrete and motivate the rest of this thesis: broadcast does not
scale gracefully, broadcast frames are unacknowledged, the distributed MAC
collides under load, and the over-water channel is harsh.

\section{IEEE 802.11 MAC fundamentals}
\subsection{Distributed Coordination Function}
The 802.11 MAC arbitrates the medium with the DCF~\cite{sharetechnote-dcf},
built on Carrier-Sense Multiple Access with Collision Avoidance
(CSMA/CA)~\cite{csma-ca}. Before transmitting, a node senses the channel; if it
is busy the node defers. Once the channel has been idle for a DCF Interframe
Space (DIFS), the node draws a random back-off counter from the contention
window $[0,CW-1]$ and decrements it slot by slot while the channel stays idle,
freezing the counter whenever the channel becomes busy and resuming when it
clears. When the counter reaches zero the node transmits.

\subsection{Contention window and the broadcast handicap}
For unicast traffic a receiver returns an acknowledgement; a missing ACK is read
as a collision and triggers a doubling of the contention window (binary
exponential back-off). \emph{Broadcast} traffic has no ACK. Consequently the
sender has no collision feedback, the contention window cannot grow in response
to load, and collisions become silent losses. In the system studied here the
contention window is therefore held fixed at $CW=16$ and reliability must be
defended by other means --- chiefly by controlling \emph{how often} nodes
choose to transmit.

\subsection{The hidden-terminal problem}
Two senders out of range of each other but both in range of a common receiver
can transmit simultaneously and collide at that receiver~\cite{tobagi-hidden}.
Unicast Wi-Fi mitigates this with the RTS/CTS handshake, which is unavailable to
broadcast. Hidden terminals are therefore a first-class source of loss in
broadcast beaconing and must be modelled faithfully (Section~\ref{sec:collmodel}).

\section{Beaconing strategies}
\subsection{Static beaconing}
The simplest scheme transmits at a fixed period regardless of conditions. It is
trivial to implement and analyze, but its channel load grows linearly with the
number of nodes; in dense scenarios it collides heavily, and in sparse ones it
wastes energy.

\subsection{Adaptive beaconing}
Adaptive schemes tune the transmission rate to local conditions. The literature,
mostly from vehicular and mobile ad-hoc networks, proposes density-based
adaptation that backs off as neighbours multiply~\cite{takai-dvcs}, mobility-based
adaptation keyed to relative velocity~\cite{alotaibi-adaptive}, and context-aware
adaptation that weighs several factors at once~\cite{sommer-adaptive}. Surveys of
beacon scheduling for device discovery confirm that adaptation can substantially
improve efficiency~\cite{chen-survey}, and VANET studies show reliability gains
from topology- and channel-aware interval control~\cite{zhao-vanet}. Most of
these designs, however, assume bidirectional links, coordinated access or IP-level
routing. The maritime setting forbids all three: communication is one-way
broadcast, there is no coordination, and there is no routing layer.

\subsection{Awareness beyond one hop}
Extending awareness past direct neighbours is classically achieved either by
\emph{piggybacking} digested neighbour information onto ordinary beacons or by
\emph{flooding} --- relaying received frames outward under a hop limit. Flooding
maximizes reach but is notorious for the ``broadcast storm''~\cite{broadcaststorm}:
in dense networks redundant rebroadcasts collide and congest the medium.
Piggybacking is cheaper but bounded in how far and how freshly it can carry
information. Which mechanism wins in a dense, unacknowledged, single-domain
broadcast network is exactly one of the questions this thesis answers
empirically.

\section{Simulation in wireless research}
Physical experiments with hundreds of nodes are expensive, hard to control and
nearly impossible to reproduce, and offshore deployments compound every
difficulty. Simulation is the standard alternative. General-purpose simulators
such as ns-3~\cite{nsnam} and OMNeT++~\cite{omnetpp} ship broad protocol stacks
but carry a steep learning curve and significant overhead when a study targets
one narrow mechanism. For focused MAC-layer questions a \emph{custom} simulator
is often preferable: it contains only what the research question needs, gives
full control over every model, and yields a compact, reproducible codebase ---
at the price of requiring careful validation, which here is provided by
field-calibrating the channel.

\section{Gap in the literature}
Wi-Fi proximity sensing has been studied indoors and in cities~\cite{goel-wifi};
maritime communication has focused on long range; adaptive beaconing has been
explored for vehicles. No prior work combines the specific constraints of
coastal proximity detection --- infrastructure-free operation, one-way
unacknowledged broadcast, energy-constrained wearable nodes, and a need for
reliable short-range \emph{and} multi-hop awareness in a churning, mobile
population --- nor does it jointly tune beaconing rate and awareness propagation
under those constraints. That joint treatment is the gap this thesis addresses.

% ================================================================
\chapter{The ProSafe System and Beaconing Protocols}
\label{ch:system}

\section{System overview}
ProSafe (Proximity-based Safety) is a minimalist, Wi-Fi-based safety network in
which \emph{smart buoys} broadcast periodic, location-aware beacons so that
vessels can detect people in the water. It runs entirely within a local ad-hoc
802.11 domain: alerts are produced and consumed locally, with no access point,
no router and no Internet. The design follows five principles --- it is
infrastructure-free, it uses pure MAC-layer broadcast (no IP), vessels detect
passively without transmitting, the hardware is low-cost commodity Wi-Fi, and
the battery-powered buoys must be energy-frugal.

\subsection{Components}
A \emph{smart user buoy} is a floating device carried by a person in the water.
It integrates a single-board computer with a Wi-Fi radio (e.g.\ a Raspberry
Pi~\cite{raspberrypi}), a GPS module, an inertial unit, a battery and a
waterproof enclosure. Its single radio both transmits its own beacons and
listens for those of others. A \emph{vessel safety unit} is installed on a boat
and listens only; when it hears a buoy above a configurable signal-strength
threshold it raises an audible and visual alert and can plot detected buoys on a
map. Because each beacon also carries the sender's recent neighbour list, the
vessel can infer the presence of buoys it has not heard directly.

\section{Communication model}
\subsection{Beacon frame}
Beacons are short broadcast frames sent to the broadcast address. The on-air
payload modelled in the simulator (class \texttt{Beacon}) is:
\begin{itemize}[leftmargin=1.4em,nosep]
  \item \textbf{sender id} --- a $16$-byte UUID;
  \item \textbf{mobile flag} --- $1$ byte;
  \item \textbf{position} --- $8$ bytes;
  \item \textbf{timestamp} --- $4$ bytes;
  \item \textbf{neighbour list} --- $28$ bytes per entry (UUID $+$ timestamp $+$
  position); the per-entry hop distance is bookkeeping and is \emph{not} counted
  on air;
  \item \textbf{origin id} and \textbf{hop limit} --- $16+4$ bytes, present only
  in forwarded mode.
\end{itemize}
The base frame is therefore $33$ bytes and grows by $28$ bytes per advertised
neighbour, so beacon size is itself a function of local density --- a fact that
matters when comparing append against forwarded mode.

\subsection{MAC operation and topology}
Channel access is standard DCF: sense, wait DIFS ($\SI{50}{\micro\second}$ in
the model), draw a back-off from $[0,CW-1]$ with $CW=16$ and
$\text{slot}=\SI{20}{\micro\second}$, then transmit. No RTS/CTS and no ACK are
used. Buoys form a pure peer-to-peer ad-hoc domain with no persistent
associations; nodes enter and leave range continuously because of currents,
waves and mobility, so each node keeps only a soft-state neighbour table that
expires stale entries.

\section{Scheduler 1: Static Beaconing Protocol (SBP)}
SBP is the baseline. Each buoy transmits every $BI_{\text{static}}$ seconds
irrespective of conditions, listens continuously, refreshes its neighbour table
on every reception, and expires entries older than a timeout. SBP is simple and
predictable but, as Chapter~\ref{ch:results} shows, its channel load and
collision rate grow with density.

\section{Scheduler 2: Adaptive Density-Aware Beaconing (ADAB)}
ADAB regulates the beacon interval using a single local quantity: the number of
neighbours $N_t$ currently in the table. It maps the count to a normalized
density factor
\begin{equation}
  f_t = \min\!\left(1,\; \frac{N_t}{\nthr}\right),
\end{equation}
and expands the interval quadratically between bounds, with a small
multiplicative jitter $\epsilon_t\sim\mathcal{U}[-\eta,\eta]$, $\eta=0.05$, to
desynchronize nodes:
\begin{equation}
  BI_{t+1} = \Big[\,\bimin + f_t^{\,2}\,(\bimax-\bimin)\,\Big]\,(1+\epsilon_t),
  \qquad BI_{t+1}\in[\bimin,\bimax].
  \label{eq:adab}
\end{equation}
With few neighbours $f_t\approx 0$ and the buoy beacons near $\bimin$ for fast
discovery; as $N_t\to\nthr$ the quadratic term drives the interval toward
$\bimax$, throttling transmissions where contention is highest. In the
implementation $\nthr=15$. ADAB changes only \emph{when} beacons are generated;
the underlying CSMA/CA is untouched.

\section{Scheduler 3: Adaptive Context-Aware Beaconing (ACAB)}
\label{sec:acab}
ADAB reacts to density alone. But density is not the only locally observable
predictor of whether a beacon is worth sending. ACAB, the scheduler introduced
in this thesis, blends three signals.

\paragraph{Density score.} As in ADAB, a normalized neighbour count
\begin{equation}
  s_{\mathrm{den}} = \min\!\left(1,\;\frac{N_t}{\nthr^{\mathrm{ACAB}}}\right),
  \qquad \nthr^{\mathrm{ACAB}} = 10 .
\end{equation}

\paragraph{Contact-recency score.} If the buoy has heard \emph{someone} recently
there is less urgency to advertise itself; if it has been silent for a while it
should speak up. With $\Delta = t - t_{\text{last contact}}$ and a recency
horizon $\tau_c = \SI{20}{\second}$,
\begin{equation}
  s_{\mathrm{con}} =
  \begin{cases}
    \max\!\big(0,\; 1 - \Delta/\tau_c\big) & N_t>0 \text{ and contact seen},\\
    0 & \text{otherwise.}
  \end{cases}
\end{equation}
A recently-contacted node yields $s_{\mathrm{con}}\to 1$ (throttle); a long-silent
or isolated node yields $0$ (beacon aggressively).

\paragraph{Mobility score.} A fast-moving node changes the picture quickly and
should refresh its position more often. With speed $v=\lVert\mathbf{v}\rVert$ and
a reference velocity $v_{\text{ref}}$,
\begin{equation}
  s_{\mathrm{mob}} = \min\!\left(1,\;\frac{v}{v_{\text{ref}}}\right).
\end{equation}

\paragraph{Blend and mapping.} The three are combined with weights
$(w_d,w_c,w_m)=(0.4,0.3,0.3)$ that sum to one, noting that mobility \emph{lowers}
the back-off (hence $1-s_{\mathrm{mob}}$):
\begin{equation}
  s = w_d\,s_{\mathrm{den}} + w_c\,s_{\mathrm{con}} + w_m\,(1 - s_{\mathrm{mob}}).
  \label{eq:acab-blend}
\end{equation}
The blended score is then fed through the same quadratic interval map and jitter
as ADAB (Eq.~\ref{eq:adab}) with $f_t\equiv s$. ACAB thus throttles a buoy that
is simultaneously dense, recently-contacted and slow --- exactly the situation in
which one more beacon adds the least value --- while keeping an isolated,
silent or fast buoy talkative. Algorithm~\ref{alg:interval} gives the unified
interval computation; the scheduler type selects which score drives $f_t$.

\begin{algorithm}[t]
\caption{Adaptive interval computation (shared by ADAB and ACAB)}
\label{alg:interval}
\begin{algorithmic}[1]
\Function{ComputeInterval}{$\mathbf{v}, N_t, t_{\text{last}}, t$}
  \If{scheduler $=$ \textsc{adab}}
     \State $s \gets \min(1, N_t / \nthr^{\mathrm{ADAB}})$
  \ElsIf{scheduler $=$ \textsc{acab}}
     \State $s_{\mathrm{den}} \gets \min(1, N_t / \nthr^{\mathrm{ACAB}})$
     \State $s_{\mathrm{con}} \gets (N_t>0 \wedge t_{\text{last}}\!\neq\!\mathrm{nil})\,?\,\max(0,1-(t-t_{\text{last}})/\tau_c) : 0$
     \State $s_{\mathrm{mob}} \gets \min(1, \lVert\mathbf{v}\rVert / v_{\text{ref}})$
     \State $s \gets w_d s_{\mathrm{den}} + w_c s_{\mathrm{con}} + w_m (1-s_{\mathrm{mob}})$
  \EndIf
  \State $f \gets s^2$ \Comment{quadratic back-off}
  \State $BI \gets \bimin + f\,(\bimax-\bimin)$
  \State $BI \gets BI \cdot (1 + \mathcal{U}[-\eta,\eta])$ \Comment{anti-sync jitter}
  \State \Return $\mathrm{clamp}(BI, \bimin, \bimax)$
\EndFunction
\end{algorithmic}
\end{algorithm}

\section{Multi-hop awareness}
\label{sec:multihop}
A beacon natively conveys one-hop awareness. ProSafe supports two ways to extend
it, selectable per run.

\subsection{Append mode (passive)}
In append mode a buoy includes, in every beacon, the nodes it has learned about
--- both its direct neighbours (hop distance $1$) and nodes discovered
indirectly through others' neighbour lists. On reception a buoy folds the
sender's advertised list into a \emph{discovered-nodes} table, incrementing each
entry's hop distance by one and dropping anything beyond an
\emph{append hop limit} (so information cannot circulate indefinitely). No extra
frame is ever sent: awareness rides for free on beacons that would be
transmitted anyway. The cost is purely in beacon \emph{size} --- $28$ bytes per
advertised node --- and hence in airtime per frame.

\subsection{Forwarded mode (active) and the forwarding gate}
In forwarded mode a buoy actively re-broadcasts beacons it receives. Each beacon
carries an \emph{origin id} and a \emph{hop limit} (TTL). On receiving a beacon
whose origin is not itself and whose TTL is positive, a buoy may relay a copy
with the TTL decremented; the origin, timestamp and payload are preserved so
duplicates can be detected and a node never re-forwards its own beacon. Relays
are paced through the same CSMA pipeline as ordinary beacons, after a small
random jitter that desynchronizes the many receivers of one frame.

Left unchecked, this is a broadcast storm waiting to happen. The thesis
introduces a \textbf{density-aware forwarding gate} that decides, once, at
enqueue time, whether a received beacon is worth relaying at all:
\begin{equation}
  P_{\text{fwd}} =
  \begin{cases}
    1 & N \le n_0,\\[2pt]
    \dfrac{n_0}{N}\cdot \max\!\big(0.3,\;1-f_q\big) & N > n_0,
  \end{cases}
  \label{eq:fwdgate}
\end{equation}
where $N$ is the relay's neighbour count, $n_0$ is a baseline expected number of
forwarders per cascade, and $f_q$ is \emph{the squared back-off score the
scheduler last computed} (line~10 of Algorithm~\ref{alg:interval}). The first
factor $n_0/N$ thins relays in dense neighbourhoods so that, in expectation,
only $\approx n_0$ of the $N$ candidates rebroadcast. The second factor is the
novel coupling: a buoy that is \emph{already} throttling its own beacons because
it sits in a busy, stable neighbourhood ($f_q$ large) further reduces its
willingness to relay, down to a floor of $0.3$. The same local signal that
governs ``how often I talk'' is thus re-used to govern ``how much I relay,''
with no new state. Suppressed beacons are recorded as seen (so later copies are
deduplicated) but never consume a contention.

\section{Energy model}
\label{sec:energy}
The dominant energy cost is radio transmission, so a buoy's hourly transmission
energy scales with its beacon rate $r$ and the per-beacon airtime
$t_{\text{air}}$. For a $500$-byte beacon at $1$\,Mbps the airtime is dominated
by the $\SI{4}{\milli\second}$ payload, giving roughly $E_{\text{beacon}}\approx
\SI{4.4}{\milli\joule}$ at a $\sim\!1$\,W transmit draw and therefore
\begin{equation}
  E_{\text{h}} = 3600\cdot r\cdot E_{\text{beacon}}.
\end{equation}
Because every adaptive decision lowers $r$ (ADAB, ACAB) or suppresses a relay
(the forwarding gate), the transmission-count metric reported in
Chapter~\ref{ch:results} is, to first order, a direct proxy for energy: halving
the beacons halves the transmission energy. ACAB's aggressive throttling and the
forwarding gate's relay suppression are, in this sense, primarily
\emph{energy} mechanisms whose reliability is shown to be preserved rather than
sacrificed.

% ================================================================
\chapter{Simulator Design and Implementation}
\label{ch:simulator}

\section{Why a custom simulator}
General-purpose simulators would force this study to navigate large protocol
stacks irrelevant to a single MAC-layer question. A custom, discrete-event
simulator instead contains exactly the buoy behaviour, channel model, beaconing
logic and metrics the work needs; it is small enough to read in full, easy to
extend with a new scheduler or awareness mode, and reproducible because every
parameter is externalized. The simulator is written in Python~\cite{python} and
organized into focused modules: \texttt{simulator.py} (engine and event queue),
\texttt{events.py}, \texttt{channel.py}, \texttt{buoy.py}, \texttt{beacon.py},
\texttt{scheduler.py}, \texttt{metrics.py} and \texttt{config\_handler.py}.

\section{Event-driven core}
The engine advances by popping events from a time-ordered priority queue
implemented with Python's \texttt{heapq}~\cite{heapq}. Each event carries its
firing time, a type, a target object (a buoy, the channel or the simulator) and
an optional payload. A monotone counter breaks ties so that events with equal
timestamps fire in insertion order, making runs deterministic for a fixed seed.
The event types span the protocol's full lifecycle: \texttt{SCHEDULER\_CHECK},
\texttt{CHANNEL\_SENSE}, \texttt{DIFS\_COMPLETION}, \texttt{BACKOFF\_COMPLETION},
\texttt{TRANSMISSION\_START}, \texttt{TRANSMISSION\_END}, \texttt{RECEPTION},
\texttt{NEIGHBOR\_CLEANUP}, \texttt{BUOY\_MOVEMENT}, \texttt{BUOY\_ARRAY\_UPDATE}
and \texttt{AVG\_NEIGHBORS\_CALCULATION}. Recurring events scheduled before a
buoy is deactivated are invalidated by a per-buoy generation counter, giving
cheap lazy cancellation when the population churns.

\section{The buoy: state and state machine}
Each buoy holds a UUID, a position, a mobility flag and velocity, a battery
level, a soft-state neighbour table, a scheduler instance, and the bookkeeping
for whichever awareness mode is active (a discovered-nodes table for append, a
pending-forward queue and a seen-origins map for forwarded). Its transmission
behaviour follows the DCF state machine in Figure~\ref{fig:fsm-desc}:

\begin{figure}[t]
\centering
\fbox{\parbox{0.92\textwidth}{\small
\textsc{receiving} $\xrightarrow{\text{scheduler decides to send}}$
\textsc{channel\_sense}:
if busy, re-sense when the medium clears;
if idle $\rightarrow$ \textsc{waiting\_difs} (wait DIFS).
\textsc{difs\_completion}: if still idle, draw a back-off and enter
\textsc{backoff}; if busy, return to sensing.
\textsc{backoff}: count down; if the channel goes busy, re-enter the full
sense--DIFS--back-off chain rather than transmitting blind.
\textsc{backoff\_completion} $\rightarrow$ \textsc{transmission\_start}: broadcast
one frame, then return to \textsc{receiving}.
}}
\caption{Textual rendering of the buoy CSMA/CA transmission state machine. The
back-off counter freezes whenever the channel becomes busy and the chain
restarts, faithfully reproducing 802.11 deferral.}
\label{fig:fsm-desc}
\end{figure}

\subsection{Interleaved transmission pipeline}
A buoy must serve two kinds of traffic: its \emph{own} scheduled beacons and, in
forwarded mode, \emph{relayed} beacons. The implementation uses a single CSMA
pipeline guarded by a \texttt{processing} lock, with a strict priority rule:
when a contention is won, an own beacon is sent if one is due, otherwise a single
queued relay is sent; if anything remains the pipeline immediately re-enters
contention. This guarantees that relay traffic can never starve a buoy's own
safety advertisement, while still draining the forward queue one frame per
contention so that relaying is paced rather than bursty. The scheduler decision
is recorded even while a relay pipeline is draining, so an own beacon is never
missed.

\section{Mobility: Random Waypoint}
Earlier work modelled mobility as straight-line motion with boundary reflection.
This thesis adopts the more realistic Random-Waypoint (RWP) model. A mobile buoy
picks a uniformly random waypoint in the arena and a speed drawn uniformly from
$[v_{\min},v_{\max}]$, advances toward the waypoint in fixed time steps, and on
arrival pauses for a uniform random time before choosing a new waypoint and
speed. RWP produces the start-stop, direction-changing tracks typical of people
moving in water far better than constant-velocity bouncing, and --- importantly
for ACAB --- it produces a realistic distribution of instantaneous speeds for
the mobility score of Eq.~\ref{eq:acab-blend}. A configurable fraction of buoys
(here $40\%$) is mobile; the rest are anchored.

\section{Channel model}
\label{sec:collmodel}
\subsection{Range-based probabilistic reception}
When a buoy broadcasts, the channel finds every active in-range receiver and
schedules a \texttt{RECEPTION} for each. Reception is probabilistic and
calibrated from over-water field measurements: with $d$ the sender--receiver
distance,
\begin{equation}
  P_{\text{rx}}(d) =
  \begin{cases}
    0.9 & d \le \SI{70}{\metre} \quad(\text{high-probability zone}),\\
    0.15 & \SI{70}{\metre} < d \le \SI{80}{\metre},\\
    0 & d > \SI{80}{\metre}.
  \end{cases}
\end{equation}
Accepted receptions are delayed by the propagation time $d/c$. An \emph{ideal}
channel mode ($P_{\text{rx}}=1$ within range) is also available to isolate MAC
collisions from path loss. This path-loss structure is decisive for interpreting
the results: because even a collision-free link delivers with probability well
below one, per-receiver delivery ratio is bounded by propagation, and the
benefit of scheduling surfaces elsewhere (Chapter~\ref{ch:results}).

\subsection{Collision and hidden-terminal modelling}
The channel tracks active transmissions as (beacon, start, end) tuples. A new
transmission collides with an existing one at a receiver $R$ when their time
windows overlap \emph{and} $R$ lies within range of the other sender, or when
the two senders are themselves in range. Crucially, a receiver in range of an
already-active sender loses the new frame even if the two senders cannot hear
each other --- this is precisely the hidden-terminal case, modelled directly.
When a later transmission collides with an earlier one at a receiver that had
already been scheduled to receive the earlier frame, that earlier reception is
\emph{revoked} (``poisoned''), reproducing the corruption a real radio would
see. Carrier sensing is modelled with a propagating wavefront: a position senses
the medium busy only once a transmission's signal, travelling at $c$, has
actually reached it and before it has passed.

\section{Metric suite}
The metrics module records, per run, a set of indicators that --- unlike a single
delivery number --- separate the distinct effects of scheduling and awareness
propagation.

\begin{description}[leftmargin=1.6em,style=nextline]
  \item[Per-receiver delivery ratio (PDR)] $=\dfrac{\text{actually received}}{\text{potentially sent}}$,
  counted at the granularity of (transmission $\times$ in-range receiver). It
  measures \emph{link-level} reliability and is sensitive to both collisions and
  path loss.
  \item[Network delivery ratio] $=\dfrac{\text{unique beacons received}}{\text{originated beacons}}$,
  counting each generated beacon once, the first time it reaches anyone, and
  \emph{excluding} forwarded copies from the denominator so the figure is
  comparable across modes. It measures whether a node's information reaches the
  network at all.
  \item[Collision rate] $=\dfrac{\text{collided receptions}}{\text{potentially sent}}$,
  a direct read-out of MAC contention.
  \item[Network-discovery percentage] --- the average fraction of all other nodes
  that a buoy has discovered (directly or via multi-hop awareness), the headline
  \emph{awareness} metric.
  \item[Reaction latency] --- the time between a beacon's generation and the first
  moment its origin is discovered by a given receiver; the safety-relevant
  ``time to first detection.''
  \item[Channel load] --- beacons sent (and, separately, forwards sent), the
  airtime/energy proxy of the energy model.
\end{description}

\section{Dynamic population (churn)}
To mimic users entering and leaving the water, the population churns. After an
initial settling delay, at random $15$--$20$\,s intervals the simulator either
removes up to a fixed fraction of the active buoys (respecting a floor of
$\max(3,\,0.2\,N_{\text{total}})$) or activates some currently-inactive ones (up
to the total). Newly activated buoys have their recurring events rescheduled.
This keeps density genuinely time-varying within a run, exercising each
scheduler's responsiveness rather than testing it only at a static operating
point.

\section{Configuration and reproducibility}
Every parameter --- world size, densities, intervals, channel calibration, CSMA
timings, scheduler thresholds and weights, multihop mode and limits --- lives in
a single YAML file read by a singleton \texttt{ConfigHandler}. A master seed is
derived once per batch and every individual simulation receives its own derived,
reproducible seed, so an entire campaign can be replayed bit-for-bit. The batch
runner sweeps densities, intervals and multihop modes in parallel, and a
companion script averages repeated runs and emits the plots used in the next
chapter.

% ================================================================
\chapter{Experimental Results}
\label{ch:results}

\section{Experimental setup}
All results below come from the field-calibrated probabilistic channel in a
$\SI{500}{\metre}\times\SI{500}{\metre}$ arena with a maximum range of
$\SI{80}{\metre}$. Each configuration was simulated for $\SI{60}{\second}$ and
repeated $30$ times with independent seeds; every reported point is the mean of
those $30$ runs and error bars, where shown, are $\pm1$ standard deviation. Node
counts ("density") range from $20$ to $100$ in steps of $20$; $40\%$ of nodes
are mobile under Random Waypoint; the population churns throughout each run. The
three schedulers (SBP, ADAB, ACAB) are evaluated under both awareness modes
(append, forwarded) and three beacon \emph{minimum} intervals
($0.25$, $0.5$, $1.0$\,s; for SBP this is the fixed interval). Unless stated
otherwise, plots use the $\SI{0.25}{\second}$ interval, where contention is
highest and differences are clearest. The CSMA parameters are
$CW=16$, $\text{DIFS}=\SI{50}{\micro\second}$,
$\text{slot}=\SI{20}{\micro\second}$, with no RTS/CTS or ACK.

Table~\ref{tab:headline} gives the scheduler comparison under append mode at the
$\SI{0.25}{\second}$ interval, which anchors the discussion that follows.

\begin{table}[t]
\centering
\caption{Scheduler comparison under \textbf{append} mode, $\SI{0.25}{\second}$
interval, averaged over $30$ runs. PDR is per-receiver delivery ratio; Net.\
disc.\ is the average percentage of the network each buoy discovers; "Sent" is
the channel-load / energy proxy.}
\label{tab:headline}
\begin{tabular}{lrrrrrr}
\toprule
Scheduler & Density & PDR & Net.\ deliv. & Coll.\ rate & Sent & Net.\ disc.\ (\%) \\
\midrule
SBP   & 20  & 0.723 & 0.681 & 0.0093 & 4\,622  & 53.9 \\
SBP   & 100 & 0.621 & 0.977 & 0.1429 & 23\,290 & 72.6 \\
ADAB  & 20  & 0.715 & 0.589 & 0.0021 & 2\,475  & 51.7 \\
ADAB  & 100 & 0.715 & 0.958 & 0.0105 & 4\,160  & 64.5 \\
ACAB  & 20  & 0.696 & 0.475 & 0.0010 & 751     & 45.0 \\
ACAB  & 100 & 0.726 & 0.971 & 0.0047 & 1\,578  & 59.3 \\
\bottomrule
\end{tabular}
\end{table}

\section{RQ1 \& RQ2: density, reliability and the cost of beaconing}
\subsection{Collisions: static degrades, adaptation stays flat}
Figure~\ref{fig:sched-collision} is the clearest single result. As density
climbs from $20$ to $100$ nodes the static scheduler's collision rate rises
super-linearly from under $1\%$ to $14.3\%$, tracking the growth in average
neighbours. Both adaptive schedulers flatten the curve almost completely: ADAB
stays at or below $1.1\%$ and ACAB below $0.5\%$ across the entire range. At the
densest operating point ACAB's collision rate is roughly $30\times$ lower than
static's and ADAB's is $14\times$ lower. This directly answers RQ1 (density
sharply degrades unmanaged broadcast) and the reliability half of RQ2
(local adaptation contains contention).

\begin{figure}[t]
\centering
\includegraphics[width=0.78\textwidth]{figures/sched_collision_append.pdf}
\caption{Collision rate versus density (append mode, $\SI{0.25}{\second}$
interval). Static beaconing degrades super-linearly with the neighbour count
(dotted, right axis); ADAB and ACAB stay nearly flat. Error bars are $\pm1$
standard deviation over $30$ runs.}
\label{fig:sched-collision}
\end{figure}

\subsection{Channel load and energy: the ACAB advantage}
The price static pays for its collisions is a transmission count that grows
without restraint. Figure~\ref{fig:sched-sent} (log scale) shows static sending
$23\,290$ beacons at density $100$ while ADAB sends $4\,160$ and ACAB only
$1\,578$ --- reductions of $82\%$ and $93\%$ respectively. Because transmission
count is a direct energy proxy (Section~\ref{sec:energy}), ACAB roughly \emph{decuples}
battery life relative to static at high density, and still nearly halves ADAB's
energy. The mechanism is exactly the one designed in Section~\ref{sec:acab}:
most buoys here are stationary ($60\%$) and, once a neighbourhood forms, recently
contacted, so ACAB's contact and mobility terms drive the interval toward
$\bimax$ far more decisively than ADAB's density term alone --- which, with an
average neighbour count that peaks near $7.7$ against a threshold of $10$--$15$,
never fully saturates.

\begin{figure}[t]
\centering
\includegraphics[width=0.78\textwidth]{figures/sched_sent_append.pdf}
\caption{Beacon transmissions versus density (append mode, $\SI{0.25}{\second}$
interval, log scale). ACAB transmits up to $93\%$ fewer beacons than static and
$62\%$ fewer than ADAB.}
\label{fig:sched-sent}
\end{figure}

\subsection{Reliability is preserved, not traded away}
A natural worry is that throttling so hard must cost reliability. It does not.
Figure~\ref{fig:sched-pdr} shows per-receiver PDR. Static's PDR \emph{falls} from
$0.72$ to $0.62$ as its collisions mount, whereas ACAB holds a flat
$\approx0.72$ across all densities --- ending up \emph{above} static by more than
ten points despite a fraction of the transmissions. The network delivery ratio
(Figure~\ref{fig:sched-deliv}) tells a complementary story: it rises with density
for every scheduler (more originators, more chances to be heard) and all three
converge above $0.95$ at high density, confirming that the throttled schedulers
still get each node's information out to the network. ACAB therefore does not
buy energy with reliability; it buys energy with \emph{redundancy} it did not
need.

\begin{figure}[t]
\centering
\begin{subfigure}{0.49\textwidth}
\includegraphics[width=\textwidth]{figures/sched_pdr_append.pdf}
\caption{Per-receiver PDR}
\label{fig:sched-pdr}
\end{subfigure}\hfill
\begin{subfigure}{0.49\textwidth}
\includegraphics[width=\textwidth]{figures/sched_deliv_append.pdf}
\caption{Network delivery ratio}
\label{fig:sched-deliv}
\end{subfigure}
\caption{Reliability under append mode at $\SI{0.25}{\second}$. Static's PDR
degrades with contention while the adaptive schedulers hold flat; network
delivery ratio rises with density and converges for all three.}
\label{fig:sched-reliability}
\end{figure}

\section{RQ4: per-link reliability versus network awareness}
The PDR and delivery-ratio curves together justify the metric refinement of
RQ4. Under the realistic path-loss channel, a collision-free link still delivers
with probability $0.9$ (near) or $0.15$ (far), so PDR is \emph{capped by
propagation}: a single delivery number conflates two very different effects.
Reporting PDR alone would have hidden ACAB's advantage --- its PDR is barely
distinguishable from ADAB's --- while reporting only network delivery ratio would
have hidden static's collision penalty, which the delivery ratio (a per-origin,
first-arrival measure) largely absorbs. Only by separating per-receiver PDR,
network delivery ratio, collision rate and channel load does the real trade-off
become legible: \emph{adaptation's value at this scale is overwhelmingly in
collisions and energy, not in raw delivery}. This is a methodological correction
to the common practice of summarizing broadcast reliability with one ratio.

\section{RQ3: passive append versus active forwarding}
We now fix the scheduler and vary the awareness mechanism.
Table~\ref{tab:modes} contrasts append and forwarded at density $100$.

\begin{table}[t]
\centering
\caption{Append versus forwarded awareness at density $100$,
$\SI{0.25}{\second}$ interval. Forwarded mode multiplies channel load and
latency without improving --- in fact slightly reducing --- network discovery.}
\label{tab:modes}
\begin{tabular}{llrrrrr}
\toprule
Scheduler & Mode & PDR & Sent & Forwards & Coll.\ rate & Net.\ disc.\ (\%) \\
\midrule
SBP  & append    & 0.621 & 23\,290 & 0      & 0.1429 & 72.6 \\
SBP  & forwarded & 0.572 & 80\,431 & 57\,123 & 0.2213 & 70.2 \\
ADAB & append    & 0.715 & 4\,160  & 0      & 0.0105 & 64.5 \\
ADAB & forwarded & 0.713 & 11\,198 & 7\,649  & 0.0298 & 61.4 \\
ACAB & append    & 0.726 & 1\,578  & 0      & 0.0047 & 59.3 \\
ACAB & forwarded & 0.716 & 4\,359  & 2\,978  & 0.0242 & 59.2 \\
\bottomrule
\end{tabular}
\end{table}

The result is consistent and, at first glance, counter-intuitive: \emph{active
relaying is the worse mechanism here on every axis that matters.}
Figure~\ref{fig:mode-sent} shows forwarded mode generating $3$--$3.5\times$ the
traffic of append (for static, $80\,431$ frames against $23\,290$, of which
$57\,123$ are relays). That extra traffic does not buy awareness:
Figure~\ref{fig:mode-netdisc} shows append discovering an \emph{equal or larger}
fraction of the network than forwarded for both schedulers, because append's
neighbour-list piggybacking carries multi-hop information with zero extra frames
while forwarded's relays mostly duplicate and collide. The collision rate
(Figure~\ref{fig:mode-collision}) is far higher under forwarding, and the
reaction latency (Figure~\ref{fig:mode-reaction}) is several-fold worse --- a
factor of about $5\times$ at the densest point ($\approx\!25$\,ms for
static-forwarded against $\approx\!5$\,ms for static-append) and rising to
$\approx\!16\times$ in the sparsest settings --- because relays sit in queues and
wait out extra contention and jitter. For a
single-domain, dense broadcast network the verdict on RQ3 is unambiguous:
\emph{prefer cheap passive piggybacking; active flooding is counter-productive.}

\begin{figure}[t]
\centering
\begin{subfigure}{0.49\textwidth}
\includegraphics[width=\textwidth]{figures/mode_sent.pdf}
\caption{Channel load (log scale)}
\label{fig:mode-sent}
\end{subfigure}\hfill
\begin{subfigure}{0.49\textwidth}
\includegraphics[width=\textwidth]{figures/mode_netdisc.pdf}
\caption{Network discovered}
\label{fig:mode-netdisc}
\end{subfigure}

\vspace{0.6em}
\begin{subfigure}{0.49\textwidth}
\includegraphics[width=\textwidth]{figures/mode_collision.pdf}
\caption{Collision rate}
\label{fig:mode-collision}
\end{subfigure}\hfill
\begin{subfigure}{0.49\textwidth}
\includegraphics[width=\textwidth]{figures/mode_reaction.pdf}
\caption{Reaction latency}
\label{fig:mode-reaction}
\end{subfigure}
\caption{Append (solid) versus forwarded (dashed) at $\SI{0.25}{\second}$.
Forwarding multiplies channel load, collisions and latency while failing to
improve --- even slightly reducing --- the fraction of the network discovered.}
\label{fig:modes}
\end{figure}

\subsection{The forwarding gate's effect}
Even though forwarding loses here, the density-aware gate of
Eq.~\ref{eq:fwdgate} makes it lose far less badly. Comparing relay counts at
density $100$: static forwarding emits $57\,123$ relays, but pairing forwarding
with ADAB cuts that to $7\,649$ and with ACAB to $2\,978$ --- an $87$--$95\%$
reduction. Part of this is the schedulers originating fewer beacons to begin
with, but the back-off-coupled factor $\max(0.3,1-f_q)$ in the gate compounds it:
a buoy already convinced (by density, recency and stillness) that it should stay
quiet also declines to relay. The gate thus turns forwarding from a guaranteed
broadcast storm into a merely expensive option, and it demonstrates a reusable
idea --- \emph{one local congestion signal can govern both origination and
relaying.}

\subsection{The value of multi-hop awareness, quantified}
It is worth making explicit how much awareness piggybacking actually adds. At
density $100$ the average neighbour count is $\approx 7.7$, so pure one-hop
awareness would let a buoy know about only $7.7/99\approx 7.8\%$ of the network.
Append mode lifts this to $59$--$73\%$ --- a $7.6$--$9.3\times$ magnification ---
at the cost of nothing but larger beacons. This is the strongest argument for
carrying neighbour lists at all, and it is achieved without a single dedicated
control frame.

\section{Interval sensitivity}
Figure~\ref{fig:interval} sweeps the minimum interval for ADAB under append.
Shrinking the floor from $1.0$\,s to $0.25$\,s predictably raises both the
transmission count and the collision rate, but the adaptive mapping keeps even
the most aggressive setting bounded: ADAB at $0.25$\,s still collides less than
$1.1\%$ at density $100$, where static at the same floor reaches $14\%$. The
interval is therefore best read as a discovery-latency knob --- smaller floors
find new neighbours faster --- whose contention cost adaptation absorbs. The same
qualitative ordering (ACAB $<$ ADAB $<$ static in load and collisions, with
preserved PDR) holds at every interval tested; the $0.25$\,s figures shown
elsewhere are the stress case, not an outlier.

\begin{figure}[t]
\centering
\begin{subfigure}{0.49\textwidth}
\includegraphics[width=\textwidth]{figures/interval_collision_adab.pdf}
\caption{Collision rate}
\end{subfigure}\hfill
\begin{subfigure}{0.49\textwidth}
\includegraphics[width=\textwidth]{figures/interval_sent_adab.pdf}
\caption{Beacon transmissions (log scale)}
\end{subfigure}
\caption{Interval sensitivity for ADAB (append mode). A smaller interval floor
trades energy and contention for faster discovery; adaptation keeps the
collision cost bounded even at the most aggressive setting.}
\label{fig:interval}
\end{figure}

\section{Discussion}
\subsection{Local information is sufficient}
Neither ADAB nor ACAB ever consults a coordinator, measures the channel
directly, or exchanges control traffic. They act on a neighbour count, a
last-contact timestamp and their own speed --- all observed for free as a side
effect of listening. That this suffices to flatten the collision curve and cut
energy by an order of magnitude is the central scientific message: in a dense,
infrastructure-free broadcast network, \emph{purely local context is enough} to
coordinate the medium implicitly.

\subsection{Parsimony beats redundancy}
Both headline findings --- ACAB beating ADAB, and append beating forwarded ---
are instances of the same principle. The system performs best when it transmits
\emph{as little as the safety task allows}. Extra beacons (static), extra
context-blind throttling headroom (ADAB versus ACAB), and extra relayed copies
(forwarded) all add channel occupancy faster than they add value once the
neighbourhood is even moderately dense. The design lesson for maritime safety
beaconing is to spend the channel only where a node is isolated, silent or
moving --- precisely where a missed beacon is most dangerous.

\subsection{Trade-offs and limits}
The savings are not free of trade-offs. ACAB discovers a somewhat smaller
fraction of the network than static at equal density ($59\%$ versus $73\%$ at
density $100$), because it simply emits less information; for a safety system
this is acceptable while the maximum interval stays bounded --- at $\bimax=5$\,s a
vessel at $10$\,m/s moves only $50$\,m, well within detection range --- but it
marks the frontier of how hard one should throttle. Likewise, the case against
forwarding is specific to a well-connected single domain; in a genuinely
partitioned topology, where two-hop neighbours are physically unreachable by
piggybacking, active relaying would regain its purpose, and the forwarding gate
exists precisely to make that regime affordable.

% ================================================================
\chapter{Conclusions}
\label{ch:conclusions}

\section{Summary of contributions}
This thesis extended a discrete-event simulator for infrastructure-free Wi-Fi
broadcast beaconing and used it to study, jointly, how often a maritime safety
buoy should beacon and how far the awareness it produces should travel. Its
contributions are a context-aware scheduler (ACAB) that fuses density, contact
recency and mobility; two awareness-propagation mechanisms (passive append and
active forwarded) together with a density-aware forwarding gate that re-uses the
scheduler's own back-off intensity; a metric suite that deliberately separates
per-link reliability from network awareness, collisions and energy; and a
field-calibrated experimental campaign over a churning, mobile, Random-Waypoint
population.

\section{Scientific insights}
Three findings stand out.
\begin{enumerate}[leftmargin=1.4em]
  \item \textbf{Context-aware adaptation is essentially free reliability and
  large energy savings.} ACAB held collisions below $0.5\%$ where static reached
  $14\%$, transmitted up to $93\%$ fewer beacons than static and $62\%$ fewer
  than ADAB, and did so with a per-receiver PDR equal to or better than either
  baseline.
  \item \textbf{Per-link reliability and network awareness must be measured
  apart.} Under a realistic path-loss channel, PDR is propagation-bounded and
  hides the value of scheduling; the benefit of adaptation lives in collision
  rate and channel airtime.
  \item \textbf{Cheap passive awareness beats expensive active flooding.} Append
  mode matched or exceeded forwarded mode's network discovery while generating
  $3$--$5\times$ less traffic and up to an order of magnitude lower reaction latency;
  multi-hop piggybacking magnified one-hop awareness roughly $8$-fold for the
  cost of larger beacons alone.
\end{enumerate}
Underlying all three is one principle: in infrastructure-free broadcast,
\emph{local information and channel parsimony} are sufficient, and usually
optimal.

\section{Limitations}
The physical layer is deliberately simplified: a piecewise-constant reception
probability, a single $1$\,Mbps rate, no capture effect and no explicit fading.
The mobility model, though improved to Random Waypoint, does not capture
correlated drift from currents or clustering of users. The scale tested
($20$--$100$ nodes in a $500\times500$\,m arena) covers realistic coastal
densities but not extreme crowding, and the energy analysis accounts for
transmission energy only, omitting reception and idle-listening draw. ACAB's
weights $(0.4,0.3,0.3)$ were fixed rather than tuned, and the append/forwarded
comparison was made in a well-connected domain.

\section{Future work}
Several directions follow naturally. A packet-length-dependent loss model would
let the study weigh append's growing beacons against forwarded's extra frames on
equal footing, and would expose whether very large neighbour lists eventually
hurt delivery. ACAB's weights and the forwarding gate's baseline invite
systematic tuning, perhaps online. Evaluating both awareness modes in
deliberately \emph{partitioned} topologies would test the hypothesis that
forwarding earns its cost only when piggybacking physically cannot reach. Real
GPS traces of swimmers and small craft would sharpen the mobility model, and a
complete energy budget including the receive chain would refine the battery
claims. Finally, the simulator and its findings transfer beyond the maritime
case to any infrastructure-free broadcast discovery problem --- vehicular safety
beacons, disaster-response ad-hoc networks, wildlife tracking and dense IoT.

\section{Final remarks}
A safe coastal network does not need a tower, a satellite, a SIM card or a
central scheduler. It needs each buoy to listen, to speak only when speaking
helps, and to share what it hears for free. This thesis has shown, on
field-calibrated simulations, that those three habits --- realized as
context-aware scheduling, passive multi-hop awareness and congestion-coupled
relaying --- turn ordinary Wi-Fi broadcast into a beacon that scales, saves
energy, and keeps the people in the water visible.

% ================================================================
\begin{thebibliography}{99}
\addcontentsline{toc}{chapter}{Bibliography}

\bibitem{imo-ais} International Maritime Organization. \emph{AIS Transponders}.
\url{https://www.imo.org/en/OurWork/Safety/Pages/AIS.aspx}, accessed 2025.

\bibitem{noaa-beacons} NOAA SARSAT. \emph{406\,MHz Emergency Distress Beacons}.
\url{https://www.sarsat.noaa.gov/emergency-406-beacons/}, accessed 2025.

\bibitem{ieee80211} IEEE 802.11 Working Group. \emph{The Working Group Setting
the Standards for Wireless LANs}. \url{https://www.ieee802.org/11/}, accessed 2025.

\bibitem{sharetechnote-dcf} ShareTechnote. \emph{WLAN DCF}.
\url{https://www.sharetechnote.com/html/WLAN_DCF.html}, accessed 2025.

\bibitem{csma-ca} Network Academy. \emph{Collision Avoidance (CSMA/CA)}.
\url{https://www.networkacademy.io/ccna/wireless/collision-avoidance-csma-ca},
accessed 2025.

\bibitem{tobagi-hidden} F.~Tobagi and L.~Kleinrock. Packet switching in radio
channels: Part II --- the hidden terminal problem in carrier sense
multiple-access and the busy-tone solution. \emph{IEEE Transactions on
Communications}, 23(12):1417--1433, 1975.

\bibitem{ipcisco-ibss} IpCisco. \emph{Wireless Principles: BSS, ESS, IBSS}.
\url{https://ipcisco.com/lesson/wireless-principles/}, accessed 2025.

\bibitem{takai-dvcs} M.~Takai, J.~Martin, R.~Bagrodia, and A.~Ren. Directional
virtual carrier sensing for directional antennas in mobile ad hoc networks. In
\emph{Proc. ACM MobiHoc}, pages 183--193, 2002.

\bibitem{alotaibi-adaptive} M.~M.~Alotaibi and H.~T.~Mouftah. Adaptive
expiration time for dynamic beacon scheduling in vehicular ad-hoc networks. In
\emph{IEEE VTC2015-Fall}, pages 1--6, 2015.

\bibitem{sommer-adaptive} C.~Sommer, O.~K.~Tonguz, and F.~Dressler. Adaptive
beaconing for delay-sensitive and congestion-aware traffic information systems.
In \emph{IEEE Vehicular Networking Conference}, pages 1--8, 2010.

\bibitem{chen-survey} L.~Chen and W.~Zhuang. Beacon scheduling for device
discovery in wireless networks: a survey. \emph{IEEE Communications Surveys \&
Tutorials}, 20(1):137--156, 2018.

\bibitem{zhao-vanet} J.~Zhao, M.~Li, Y.~Wang, and J.~Wang. An adaptive beaconing
strategy based on local topology and channel status in VANETs. \emph{IEEE
Transactions on Mobile Computing}, 18(10):2354--2367, 2019.

\bibitem{broadcaststorm} S.-Y.~Ni, Y.-C.~Tseng, Y.-S.~Chen, and J.-P.~Sheu. The
broadcast storm problem in a mobile ad hoc network. In \emph{Proc. ACM
MobiCom}, pages 151--162, 1999.

\bibitem{goel-wifi} M.~Goel and S.~N.~Patel. Wi-Fi-based proximity detection in
the wild. In \emph{Proc. ACM MobiSys}, pages 1--12, 2020.

\bibitem{nsnam} nsnam.org. \emph{ns-3 Network Simulator}.
\url{https://www.nsnam.org/}, accessed 2025.

\bibitem{omnetpp} omnetpp.org. \emph{OMNeT++ Discrete Event Simulator}.
\url{https://omnetpp.org/}, accessed 2025.

\bibitem{raspberrypi} Raspberry Pi Ltd. \emph{Raspberry Pi}.
\url{https://www.raspberrypi.com/}, accessed 2025.

\bibitem{python} Python Software Foundation. \emph{Python}.
\url{https://www.python.org/}, accessed 2025.

\bibitem{heapq} Python docs. \emph{heapq --- Heap queue algorithm}.
\url{https://docs.python.org/3/library/heapq.html}, accessed 2025.

\end{thebibliography}

\end{document}
